\clearpage
\phantomsection
\addcontentsline{toc}{chapter}{CHƯƠNG 3: THIẾT KẾ LOGIC — PHỤ THUỘC HÀM VÀ CHUẨN HÓA}
{\LARGE\bfseries CHƯƠNG 3: THIẾT KẾ LOGIC — PHỤ THUỘC HÀM VÀ CHUẨN HÓA}\par\vspace{1.2em}

Chương này minh chứng việc áp dụng lý thuyết phụ thuộc hàm, bao đóng, tìm khóa và các thuật toán chuẩn hóa (Chương 3) để suy ra đúng lược đồ quan hệ đã dùng cho nhóm Sự kiện của FamilyConnect, sau đó tổng hợp kết quả chuẩn hóa cho toàn bộ hệ thống.

\vspace{1.8em}
\phantomsection
\addcontentsline{toc}{section}{3.1. Ví dụ minh chứng: từ bảng phẳng đến lược đồ chuẩn hóa}
{\Large\bfseries 3.1. Ví dụ minh chứng: từ bảng phẳng đến lược đồ chuẩn hóa}\par\vspace{0.8em}

Giả sử ban đầu module Sự kiện được thiết kế bằng một bảng báo cáo phẳng duy nhất (cách làm trực giác nhưng sai, giống ví dụ NV\_DUAN/Star Wars của Chương 3):

\begin{lstlisting}[language=SQL]
	BAO_CAO_SU_KIEN(maSuKien, tieuDe, thoiGianBD, maGiaDinh, tenGiaDinh, maThanhVien, hoTenTV, rsvpStatus)
\end{lstlisting}

\vspace{1.4em}
\phantomsection
\addcontentsline{toc}{subsection}{3.1.1. Xác định phụ thuộc hàm từ quy tắc nghiệp vụ}
{\large\bfseries 3.1.1. Xác định phụ thuộc hàm từ quy tắc nghiệp vụ}\par\vspace{0.6em}

\begin{itemize}
	\item F1: maSuKien $\rightarrow$ tieuDe, thoiGianBD, maGiaDinh (mỗi sự kiện có tiêu đề, thời gian và thuộc một dòng họ duy nhất)
	\item F2: maGiaDinh $\rightarrow$ tenGiaDinh (mỗi mã dòng họ xác định một tên dòng họ)
	\item F3: maThanhVien $\rightarrow$ hoTenTV, maGiaDinh (mỗi thành viên có họ tên và thuộc một dòng họ cố định)
	\item F4: (maSuKien, maThanhVien) $\rightarrow$ rsvpStatus (trạng thái RSVP chỉ có nghĩa trong ngữ cảnh một lần mời cụ thể)
\end{itemize}

\emph{Lưu ý: F3 là giả định đơn giản hóa chỉ dành riêng cho ví dụ minh họa phân rã ở mục này (mỗi thành viên gắn với đúng một dòng họ khai sinh/gốc). Giả định này khác với ngữ cảnh phân quyền đa dòng họ đã trình bày ở Mục 2.5 và bảng UserRoles(UserId, RoleId, FamilyId), nơi một người dùng có thể giữ vai trò ở nhiều dòng họ khác nhau — hai chỗ mô tả hai khái niệm khác nhau (dòng họ gốc của thành viên vs. dòng họ mà người dùng có quyền truy cập), không mâu thuẫn nhau.}

\vspace{1.4em}
\phantomsection
\addcontentsline{toc}{subsection}{3.1.2. Tìm khóa bằng bao đóng}
{\large\bfseries 3.1.2. Tìm khóa bằng bao đóng}\par\vspace{0.6em}

Thuộc tính rsvpStatus chỉ xuất hiện ở vế phải của F4, và không FD nào khác xác định được nó ngoài (maSuKien, maThanhVien) — do đó mọi khóa phải chứa cặp này. Kiểm tra bao đóng:

\begin{lstlisting}[language=SQL]
	{maSuKien, maThanhVien}+ :
	Buoc 0: X+ = {maSuKien, maThanhVien}
	Ap dung F1 (maSuKien <= X+): X+ = {maSuKien, maThanhVien, tieuDe, thoiGianBD, maGiaDinh}
	Ap dung F3 (maThanhVien <= X+): X+ = X+ U {hoTenTV} (maGiaDinh da co)
	Ap dung F2 (maGiaDinh <= X+): X+ = X+ U {tenGiaDinh}
	Ap dung F4: X+ = X+ U {rsvpStatus}
	=> X+ = U (toan bo thuoc tinh)
\end{lstlisting}

Vậy K = (maSuKien, maThanhVien) là siêu khóa. Vì maSuKien+ $\neq$ U và maThanhVien+ $\neq$ U (mỗi vế chỉ suy ra được một phần), K là khóa ứng viên tối thiểu duy nhất của BAO\_CAO\_SU\_KIEN.

\vspace{1.4em}
\phantomsection
\addcontentsline{toc}{subsection}{3.1.3. Bất thường dữ liệu (data anomalies)}
{\large\bfseries 3.1.3. Bất thường dữ liệu (data anomalies)}\par\vspace{0.6em}

\begingroup
\renewcommand{\arraystretch}{1.25}
\setlength{\tabcolsep}{5pt}
\begin{longtable}{|>{\raggedright\arraybackslash}p{0.25\textwidth}|>{\raggedright\arraybackslash}p{0.68\textwidth}|}
	\hline
	\textbf{Loại bất thường} & \textbf{Biểu hiện trên BAO\_CAO\_SU\_KIEN} \\ \hline
	\endfirsthead
	\hline
	\textbf{Loại bất thường} & \textbf{Biểu hiện trên BAO\_CAO\_SU\_KIEN} \\ \hline
	\endhead
	Thêm (Insertion) & Không thể lưu một sự kiện mới nếu chưa có ít nhất một thành viên được mời, vì khóa chính bắt buộc có maThanhVien. \\ \hline
	Sửa (Update) & Đổi tieuDe của một sự kiện có 50 người tham gia phải sửa 50 dòng, dễ gây mâu thuẫn dữ liệu nếu sửa sót. \\ \hline
	Xóa (Deletion) & Xóa người tham gia cuối cùng của một sự kiện sẽ xóa luôn thông tin tieuDe, thoiGianBD, maGiaDinh của sự kiện đó. \\ \hline
\end{longtable}
\endgroup

\vspace{1.4em}
\phantomsection
\addcontentsline{toc}{subsection}{3.1.4. Kiểm tra dạng chuẩn và phân rã (2NF $\rightarrow$ 3NF $\rightarrow$ BCNF)}
{\large\bfseries 3.1.4. Kiểm tra dạng chuẩn và phân rã (2NF $\rightarrow$ 3NF $\rightarrow$ BCNF)}\par\vspace{0.6em}

Với khóa K = (maSuKien, maThanhVien):

\begin{itemize}
	\item Vi phạm 2NF: tieuDe, thoiGianBD, maGiaDinh, tenGiaDinh chỉ phụ thuộc vào phần maSuKien của khóa (phụ thuộc từng phần); hoTenTV chỉ phụ thuộc vào phần maThanhVien.
	\item Áp dụng quy tắc tách 2NF (Chương 3, mục 'Phân rã để đạt 2NF'): tách các thuộc tính phụ thuộc từng phần ra khỏi bảng gốc theo từng phần khóa mà chúng phụ thuộc.
\end{itemize}

\begin{lstlisting}[language=SQL]
	Sau khi tach 2NF:
	SU_KIEN_TMP(maSuKien, tieuDe, thoiGianBD, maGiaDinh, tenGiaDinh)
	THANH_VIEN_TMP(maThanhVien, hoTenTV, maGiaDinh)
	THAM_GIA(maSuKien, maThanhVien, rsvpStatus)
\end{lstlisting}

Trong SU\_KIEN\_TMP vẫn còn phụ thuộc bắc cầu: maSuKien $\rightarrow$ maGiaDinh $\rightarrow$ tenGiaDinh, với maGiaDinh không phải là siêu khóa của SU\_KIEN\_TMP $\Rightarrow$ vi phạm 3NF. Áp dụng tiếp quy tắc tách 3NF (loại phụ thuộc bắc cầu):

\begin{lstlisting}[language=SQL]
	Sau khi tach 3NF:
	GIA_DINH(maGiaDinh, tenGiaDinh)
	SU_KIEN(maSuKien, tieuDe, thoiGianBD, maGiaDinh FK)
	THANH_VIEN(maThanhVien, hoTenTV, maGiaDinh FK)
	THAM_GIA(maSuKien FK, maThanhVien FK, rsvpStatus)
\end{lstlisting}

Bốn lược đồ trên chính là các bảng FAMILIES, EVENTS, FAMILY\_MEMBERS và EVENT\_PARTICIPANTS đã dùng trong thiết kế thực tế (Chương 2), xác nhận thiết kế trực giác ban đầu là đúng theo lý thuyết chuẩn hóa hình thức.

\vspace{1.4em}
\phantomsection
\addcontentsline{toc}{subsection}{3.1.5. Kiểm tra BCNF}
{\large\bfseries 3.1.5. Kiểm tra BCNF}\par\vspace{0.6em}

\begingroup
\renewcommand{\arraystretch}{1.25}
\setlength{\tabcolsep}{5pt}
\begin{longtable}{|>{\raggedright\arraybackslash}p{0.22\textwidth}|>{\raggedright\arraybackslash}p{0.38\textwidth}|>{\raggedright\arraybackslash}p{0.32\textwidth}|}
	\hline
	\textbf{Quan hệ} & \textbf{FD không tầm thường} & \textbf{X có là siêu khóa? $\rightarrow$ BCNF?} \\ \hline
	\endfirsthead
	\hline
	\textbf{Quan hệ} & \textbf{FD không tầm thường} & \textbf{X có là siêu khóa? $\rightarrow$ BCNF?} \\ \hline
	\endhead
	GIA\_DINH & maGiaDinh $\rightarrow$ tenGiaDinh & Có (maGiaDinh là khóa) $\rightarrow$ Đạt BCNF \\ \hline
	SU\_KIEN & maSuKien $\rightarrow$ \{tieuDe, thoiGianBD, maGiaDinh\} & Có (maSuKien là khóa) $\rightarrow$ Đạt BCNF \\ \hline
	THANH\_VIEN & maThanhVien $\rightarrow$ \{hoTenTV, maGiaDinh\} & Có (maThanhVien là khóa) $\rightarrow$ Đạt BCNF \\ \hline
	THAM\_GIA & (maSuKien, maThanhVien) $\rightarrow$ rsvpStatus & Vế trái là khóa (khóa ghép) $\rightarrow$ Đạt BCNF \\ \hline
\end{longtable}
\endgroup

\vspace{1.4em}
\phantomsection
\addcontentsline{toc}{subsection}{3.1.6. Kiểm tra bảo toàn kết nối (lossless-join)}
{\large\bfseries 3.1.6. Kiểm tra bảo toàn kết nối (lossless-join)}\par\vspace{0.6em}

Áp dụng tiêu chí giao hai lược đồ phải là khóa của ít nhất một bên (Chương 3, mục 'Bảo toàn kết nối nhìn bằng phép kết'):

\begin{itemize}
	\item SU\_KIEN $\cap$ GIA\_DINH = \{maGiaDinh\}, là khóa của GIA\_DINH $\Rightarrow$ bảo toàn kết nối.
	\item THANH\_VIEN $\cap$ GIA\_DINH = \{maGiaDinh\}, là khóa của GIA\_DINH $\Rightarrow$ bảo toàn kết nối.
	\item THAM\_GIA $\cap$ SU\_KIEN = \{maSuKien\}, là khóa của SU\_KIEN $\Rightarrow$ bảo toàn kết nối.
	\item THAM\_GIA $\cap$ THANH\_VIEN = \{maThanhVien\}, là khóa của THANH\_VIEN $\Rightarrow$ bảo toàn kết nối.
\end{itemize}

Phép phân rã cũng bảo toàn phụ thuộc: cả bốn FD gốc (F1–F4) đều được kiểm tra được cục bộ trong từng bảng con, không cần phép kết để xác nhận ràng buộc — đạt cả hai tiêu chí mong muốn của một phép phân rã tốt.

\vspace{1.8em}
\phantomsection
\addcontentsline{toc}{section}{3.2. Tổng hợp phân tích chuẩn hóa cho toàn bộ hệ thống}
{\Large\bfseries 3.2. Tổng hợp phân tích chuẩn hóa cho toàn bộ hệ thống}\par\vspace{0.8em}

Áp dụng cùng quy trình (xác định phụ thuộc hàm từ quy tắc nghiệp vụ $\rightarrow$ tìm khóa bằng bao đóng $\rightarrow$ kiểm tra dạng chuẩn 2NF/3NF/BCNF) cho các nhóm chức năng còn lại, kết quả tổng hợp được trình bày trong bảng bên dưới.

Đối với các bảng sử dụng khóa chính là khóa thay thế (UUID), tất cả thuộc tính không khóa đều phụ thuộc đầy đủ và không bắc cầu vào khóa chính. Các bảng trung gian ($n-n$) có khóa ghép trùng khớp hoàn toàn với vế trái của phụ thuộc hàm định nghĩa chúng.

\vspace{0.8em}
\begingroup
\footnotesize
\renewcommand{\arraystretch}{1.3}
\setlength{\tabcolsep}{2.5pt}
\begin{longtable}{|>{\raggedright\arraybackslash}p{0.14\textwidth}
		|>{\raggedright\arraybackslash}p{0.31\textwidth}
		|>{\raggedright\arraybackslash}p{0.30\textwidth}
		|>{\raggedright\arraybackslash}p{0.13\textwidth}
		|>{\centering\arraybackslash}p{0.07\textwidth}|}
	\hline
	\textbf{Quan hệ} & \textbf{Lược đồ quan hệ (Mô hình đầy đủ)} & \textbf{Phụ thuộc hàm đầy đủ (FD)} & \textbf{Khóa chính} & \textbf{Dạng chuẩn} \\ \hline
	\endfirsthead
	\hline
	\textbf{Quan hệ} & \textbf{Lược đồ quan hệ (Mô hình đầy đủ)} & \textbf{Phụ thuộc hàm đầy đủ (FD)} & \textbf{Khóa chính} & \textbf{Dạng chuẩn} \\ \hline
	\endhead
	
	Users & Users(\textbf{UserId}, Username, Email, Phone, PasswordHash, FullName, AvatarUrl, Status, LastLoginAt, CreatedAt, UpdatedAt) & UserId $\rightarrow$ Username, Email, Phone, PasswordHash, FullName, AvatarUrl, Status, LastLoginAt, CreatedAt, UpdatedAt & UserId & BCNF \\ \hline
	
	UserProfiles & UserProfiles(\textbf{UserId}, DateOfBirth, Gender, Address, Bio, UpdatedAt) & UserId $\rightarrow$ DateOfBirth, Gender, Address, Bio, UpdatedAt & UserId (PK/FK) & BCNF \\ \hline
	
	Roles & Roles(\textbf{RoleId}, RoleName, Description) & RoleId $\rightarrow$ RoleName, Description & RoleId & BCNF \\ \hline
	
	Permissions & Permissions(\textbf{PermissionId}, PermissionCode, Description) & PermissionId $\rightarrow$ PermissionCode, Description & PermissionId & BCNF \\ \hline
	
	RolePermissions & RolePermissions(\textbf{RoleId}, \textbf{PermissionId}) & (RoleId, PermissionId) $\rightarrow \emptyset$ & (RoleId, PermissionId) & BCNF \\ \hline
	
	Families & Families(\textbf{FamilyId}, FamilyName, Origin, Description, CoverImageUrl, CreatedBy, CreatedAt) & FamilyId $\rightarrow$ FamilyName, Origin, Description, CoverImageUrl, CreatedBy, CreatedAt & FamilyId & BCNF \\ \hline
	
	UserRoles & UserRoles(\textbf{UserId}, \textbf{RoleId}, \textbf{FamilyId}) & (UserId, RoleId, FamilyId) $\rightarrow \emptyset$ & (UserId, RoleId, FamilyId) & BCNF \\ \hline
	
	Branches & Branches(\textbf{BranchId}, FamilyId, ParentBranchId, BranchName, Description) & BranchId $\rightarrow$ FamilyId, ParentBranchId, BranchName, Description & BranchId & BCNF \\ \hline
	
	Family\allowbreak Members & FamilyMembers(\textbf{MemberId}, FamilyId, BranchId, UserId, FullName, Gender, DateOfBirth, DateOfDeath, IsAlive, BirthPlace, GenerationLevel, Biography, AvatarUrl, CreatedAt) & MemberId $\rightarrow$ FamilyId, BranchId, UserId, FullName, Gender, DateOfBirth, DateOfDeath, IsAlive, BirthPlace, GenerationLevel, Biography, AvatarUrl, CreatedAt & MemberId & BCNF \\ \hline
	
	ParentChild\allowbreak Relationships & ParentChild\allowbreak Relationships(\textbf{Relation\allowbreak shipId}, ParentMemberId, ChildMemberId, RelationshipType) & RelationshipId $\rightarrow$ ParentMemberId, ChildMemberId, RelationshipType; UNIQUE(Parent\allowbreak MemberId, ChildMemberId) & Relation\allowbreak shipId & BCNF \\ \hline
	
	Marriages & Marriages(\textbf{MarriageId}, Spouse1Id, Spouse2Id, MarriageDate, DivorceDate, Status) & MarriageId $\rightarrow$ Spouse1Id, Spouse2Id, MarriageDate, DivorceDate, Status & MarriageId & BCNF \\ \hline
	
	Member\allowbreak Verifications & Member\allowbreak Verifications(\textbf{Verifi\allowbreak cationId}, MemberId, RequestedBy, Status, EvidenceUrl, VerifiedBy, VerifiedAt, CreatedAt) & VerificationId $\rightarrow$ MemberId, RequestedBy, Status, EvidenceUrl, VerifiedBy, VerifiedAt, CreatedAt & Verifi\allowbreak cationId & BCNF \\ \hline
	
	Posts & Posts(\textbf{PostId}, FamilyId, AuthorId, PostType, Content, Visibility, IsPinned, CreatedAt, UpdatedAt) & PostId $\rightarrow$ FamilyId, AuthorId, PostType, Content, Visibility, IsPinned, CreatedAt, UpdatedAt & PostId & BCNF \\ \hline
	
	Comments & Comments(\textbf{CommentId}, PostId, AuthorId, ParentCommentId, Content, CreatedAt) & CommentId $\rightarrow$ PostId, AuthorId, ParentCommentId, Content, CreatedAt & CommentId & BCNF \\ \hline
	
	Reactions & Reactions(\textbf{ReactionId}, UserId, TargetType, TargetId, ReactionType, CreatedAt) & ReactionId $\rightarrow$ UserId, TargetType, TargetId, ReactionType, CreatedAt; UNIQUE(UserId, TargetType, TargetId) & ReactionId & BCNF \\ \hline
	
	MediaAssets & MediaAssets(\textbf{MediaId}, OwnerType, OwnerId, MediaUrl, MediaType, Caption, UploadedBy, UploadedAt) & MediaId $\rightarrow$ OwnerType, OwnerId, MediaUrl, MediaType, Caption, UploadedBy, UploadedAt & MediaId & BCNF \\ \hline
	
	Events & Events(\textbf{EventId}, FamilyId, CreatedBy, Title, Description, EventType, Location, StartTime, EndTime, CreatedAt) & EventId $\rightarrow$ FamilyId, CreatedBy, Title, Description, EventType, Location, StartTime, EndTime, CreatedAt & EventId & BCNF \\ \hline
	
	Event\allowbreak Participants & EventParticipants(\textbf{EventId}, \textbf{MemberId}, RsvpStatus, RespondedAt) & (EventId, MemberId) $\rightarrow$ RsvpStatus, RespondedAt & (EventId, MemberId) & BCNF \\ \hline
	
	EventReminders & EventReminders(\textbf{EventId}, \textbf{RemindAt}, Channel, IsSent) & (EventId, RemindAt) $\rightarrow$ Channel, IsSent & (EventId, RemindAt) & BCNF \\ \hline
	
	Professional\allowbreak Profiles & ProfessionalProfiles(\textbf{ProfileId}, MemberId, Occupation, Company, Industry, Position, YearsExperience, Location) & ProfileId $\rightarrow$ MemberId, Occupation, Company, Industry, Position, YearsExperience, Location & ProfileId & BCNF \\ \hline
	
	Education\allowbreak Profiles & EducationProfiles(\textbf{EducationId}, MemberId, SchoolName, Degree, FieldOfStudy, StartYear, EndYear) & EducationId $\rightarrow$ MemberId, SchoolName, Degree, FieldOfStudy, StartYear, EndYear & EducationId & BCNF \\ \hline
	
	Historical\allowbreak Documents & HistoricalDocuments(\textbf{DocumentId}, FamilyId, Title, Description, DocType, FileUrl, UploadedBy, UploadedAt) & DocumentId $\rightarrow$ FamilyId, Title, Description, DocType, FileUrl, UploadedBy, UploadedAt & DocumentId & BCNF \\ \hline
	
	FamilyStories & FamilyStories(\textbf{StoryId}, FamilyId, RelatedMemberId, AuthorId, Title, Content, CreatedAt) & StoryId $\rightarrow$ FamilyId, RelatedMemberId, AuthorId, Title, Content, CreatedAt & StoryId & BCNF \\ \hline
	
	Outstanding\allowbreak Members & Outstanding\allowbreak Members(\textbf{Achieve\allowbreak mentId}, MemberId, AchievementTitle, Description, AchievementYear) & AchievementId $\rightarrow$ MemberId, AchievementTitle, Description, AchievementYear & Achieve\allowbreak mentId & BCNF \\ \hline
	
	PhotoAlbums & PhotoAlbums(\textbf{AlbumId}, FamilyId, AlbumName, Description, CreatedBy, CreatedAt) & AlbumId $\rightarrow$ FamilyId, AlbumName, Description, CreatedBy, CreatedAt & AlbumId & BCNF \\ \hline
	
	AiQueryLogs & AiQueryLogs(\textbf{QueryId}, UserId, QueryType, QueryText, ResponseText, CreatedAt) & QueryId $\rightarrow$ UserId, QueryType, QueryText, ResponseText, CreatedAt & QueryId & BCNF \\ \hline
	
	AiEmbeddings & AiEmbeddings(\textbf{EmbeddingId}, EntityType, EntityId, EmbeddingVector, ModelVersion, CreatedAt) & EmbeddingId $\rightarrow$ EntityType, EntityId, EmbeddingVector, ModelVersion, CreatedAt & EmbeddingId & BCNF \\ \hline
	
	Ai\allowbreak Recom\allowbreak mendations & AiRecom\allowbreak mendations(\textbf{Recom\allowbreak mendationId}, UserId, TargetType, TargetId, Score, Reason, CreatedAt) & Recom\allowbreak mendationId $\rightarrow$ UserId, TargetType, TargetId, Score, Reason, CreatedAt & Recom\allowbreak mendationId & BCNF \\ \hline
	
	Reports & Reports(\textbf{ReportId}, FamilyId, ReportType, Parameters, FileUrl, GeneratedBy, GeneratedAt) & ReportId $\rightarrow$ FamilyId, ReportType, Parameters, FileUrl, GeneratedBy, GeneratedAt & ReportId & BCNF \\ \hline
	
	AuditLogs & AuditLogs(\textbf{LogId}, UserId, Action, EntityType, EntityId, OldValue, NewValue, IpAddress, CreatedAt) & LogId $\rightarrow$ UserId, Action, EntityType, EntityId, OldValue, NewValue, IpAddress, CreatedAt & LogId & BCNF \\ \hline
	
	System\allowbreak Configu\allowbreak rations & System\allowbreak Configurations(\textbf{ConfigId}, ConfigKey, ConfigValue, Description, UpdatedBy, UpdatedAt) & ConfigId $\rightarrow$ ConfigKey, ConfigValue, Description, UpdatedBy, UpdatedAt; UNIQUE(ConfigKey) & ConfigId & BCNF \\ \hline
	
	Backups & Backups(\textbf{BackupId}, BackupType, FilePath, Status, CreatedAt) & BackupId $\rightarrow$ BackupType, FilePath, Status, CreatedAt & BackupId & BCNF \\ \hline
	
	Moderation\allowbreak Logs & ModerationLogs(\textbf{ModerationId}, ContentType, ContentId, ModeratorId, Action, Reason, CreatedAt) & ModerationId $\rightarrow$ ContentType, ContentId, ModeratorId, Action, Reason, CreatedAt & ModerationId & BCNF \\ \hline
	
\end{longtable}
\endgroup

\vspace{0.8em}
\noindent\textbf{Nhận xét chung:} Việc sử dụng khóa thay thế (UUID) đơn cột cho phần lớn thực thể mạnh giúp mọi phụ thuộc hàm không tầm thường đều có vế trái là khóa (thỏa mãn BCNF một cách tự nhiên). Đồng thời, không tồn tại phụ thuộc hàm bắc cầu giữa các thuộc tính không khóa. Các bảng trung gian ($n-n$) có khóa ghép trùng khớp với vế trái duy nhất của phụ thuộc hàm định nghĩa dữ liệu quan hệ, do đó cũng đạt chuẩn BCNF hiển nhiên.

\vspace{1.8em}
\phantomsection
\addcontentsline{toc}{section}{3.3. Ràng buộc toàn vẹn (RBTV)}
{\Large\bfseries 3.3. Ràng buộc toàn vẹn (RBTV)}\par\vspace{0.8em}

Bên cạnh các ràng buộc khóa và khóa ngoại đã suy ra từ chuẩn hóa (Mục 3.1–3.2), FamilyConnect còn có một số quy tắc nghiệp vụ không thể biểu diễn thuần túy bằng khóa chính/khóa ngoại. Các quy tắc này được phát biểu tường minh dưới dạng Ràng buộc toàn vẹn (RBTV) kèm bảng tầm ảnh hưởng (ảnh hưởng trên các thao tác \texttt{Insert}, \texttt{Delete}, \texttt{Update}) theo phương pháp luận môn học.

\vspace{1.4em}
\phantomsection
\addcontentsline{toc}{subsection}{3.3.1. Ràng buộc toàn vẹn có bối cảnh là một quan hệ}
{\large\bfseries 3.3.1. Ràng buộc toàn vẹn có bối cảnh là một quan hệ}\par\vspace{0.6em}

\textbf{RB-1/Users (RBTV miền giá trị)}
\begin{itemize}
	\item \emph{Nội dung:} Trạng thái tài khoản chỉ nhận một trong bốn giá trị hợp lệ.
\end{itemize}

\begin{lstlisting}[language=SQL]
	forall u in Users (u.Status in {'active', 'inactive', 'suspended', 'pending'})
\end{lstlisting}
\noindent\textbf{Bối cảnh:} \texttt{Users}.

\vspace{0.4em}
{\centering
	\renewcommand{\arraystretch}{1.2}
	\begin{tabular}{|>{\raggedright\arraybackslash}p{0.22\textwidth}|>{\centering\arraybackslash}p{0.22\textwidth}|>{\centering\arraybackslash}p{0.22\textwidth}|>{\centering\arraybackslash}p{0.22\textwidth}|}
		\hline
		\textbf{RB-1} & \textbf{Insert} & \textbf{Delete} & \textbf{Update} \\ \hline
		Users & + & - & + \\ \hline
	\end{tabular}\par}

\vspace{1.0em}
\textbf{RB-2/Events (RBTV liên thuộc tính)}
\begin{itemize}
	\item \emph{Nội dung:} Thời điểm bắt đầu sự kiện phải diễn ra trước thời điểm kết thúc (nếu có).
\end{itemize}

\begin{lstlisting}[language=SQL]
	forall e in Events (e.EndTime IS NULL OR e.StartTime <= e.EndTime)
\end{lstlisting}
\noindent\textbf{Bối cảnh:} \texttt{Events}.

\vspace{0.4em}
{\centering
	\renewcommand{\arraystretch}{1.2}
	\begin{tabular}{|>{\raggedright\arraybackslash}p{0.22\textwidth}|>{\centering\arraybackslash}p{0.22\textwidth}|>{\centering\arraybackslash}p{0.22\textwidth}|>{\centering\arraybackslash}p{0.22\textwidth}|}
		\hline
		\textbf{RB-2} & \textbf{Insert} & \textbf{Delete} & \textbf{Update} \\ \hline
		Events & + & - & + \\ \hline
	\end{tabular}\par}

\vspace{1.0em}
\textbf{RB-3/FamilyMembers (RBTV liên thuộc tính)}
\begin{itemize}
	\item \emph{Nội dung:} Thành viên đã mất bắt buộc phải có ngày mất; thành viên còn sống không được có ngày mất.
\end{itemize}

\begin{lstlisting}[language=SQL]
	forall m in FamilyMembers ((m.IsAlive = 0 => m.DateOfDeath IS NOT NULL) AND (m.IsAlive = 1 => m.DateOfDeath IS NULL))
\end{lstlisting}
\noindent\textbf{Bối cảnh:} \texttt{FamilyMembers}.

\vspace{0.4em}
{\centering
	\renewcommand{\arraystretch}{1.2}
	\begin{tabular}{|>{\raggedright\arraybackslash}p{0.22\textwidth}|>{\centering\arraybackslash}p{0.22\textwidth}|>{\centering\arraybackslash}p{0.22\textwidth}|>{\centering\arraybackslash}p{0.22\textwidth}|}
		\hline
		\textbf{RB-3} & \textbf{Insert} & \textbf{Delete} & \textbf{Update} \\ \hline
		FamilyMembers & + & - & + \\ \hline
	\end{tabular}\par}

\vspace{1.0em}
\textbf{RB-4/ParentChildRelationships (RBTV liên bộ, tự tham chiếu)}
\begin{itemize}
	\item \emph{Nội dung:} Một thành viên không thể vừa là cha/mẹ vừa là con của chính mình trong cùng một quan hệ.
\end{itemize}

\begin{lstlisting}[language=SQL]
	forall r in ParentChildRelationships (r.ParentMemberId != r.ChildMemberId)
\end{lstlisting}
\noindent\textbf{Bối cảnh:} \texttt{ParentChildRelationships}. Ràng buộc này đã được hiện thực bằng \texttt{CONSTRAINT CK\_NotSelfParent} trong Phụ lục A.

\vspace{0.4em}
{\centering
	\renewcommand{\arraystretch}{1.2}
	\begin{tabular}{|>{\raggedright\arraybackslash}p{0.22\textwidth}|>{\centering\arraybackslash}p{0.22\textwidth}|>{\centering\arraybackslash}p{0.22\textwidth}|>{\centering\arraybackslash}p{0.22\textwidth}|}
		\hline
		\textbf{RB-4} & \textbf{Insert} & \textbf{Delete} & \textbf{Update} \\ \hline
		ParentChildRelationships & + & - & + \\ \hline
	\end{tabular}\par}

\vspace{1.4em}
\phantomsection
\addcontentsline{toc}{subsection}{3.3.2. Ràng buộc toàn vẹn có bối cảnh là nhiều quan hệ}
{\large\bfseries 3.3.2. Ràng buộc toàn vẹn có bối cảnh là nhiều quan hệ}\par\vspace{0.6em}

\textbf{RB-5/EventParticipants, Events, FamilyMembers (RBTV tham chiếu liên quan hệ)}
\begin{itemize}
	\item \emph{Nội dung:} Thành viên được mời tham gia sự kiện phải thuộc cùng dòng họ với dòng họ tổ chức sự kiện đó.
\end{itemize}

\begin{lstlisting}[language=SQL]
	forall p in EventParticipants, exists e in Events, exists m in FamilyMembers
	(e.EventId = p.EventId AND m.MemberId = p.MemberId AND e.FamilyId = m.FamilyId)
\end{lstlisting}
\noindent\textbf{Bối cảnh:} \texttt{EventParticipants}, \texttt{Events}, \texttt{FamilyMembers}.

\vspace{0.4em}
{\centering
	\renewcommand{\arraystretch}{1.2}
	\begin{tabular}{|>{\raggedright\arraybackslash}p{0.22\textwidth}|>{\centering\arraybackslash}p{0.22\textwidth}|>{\centering\arraybackslash}p{0.22\textwidth}|>{\centering\arraybackslash}p{0.22\textwidth}|}
		\hline
		\textbf{RB-5} & \textbf{Insert} & \textbf{Delete} & \textbf{Update} \\ \hline
		EventParticipants & + & - & + \\ \hline
		Events & - & - & + \\ \hline
		FamilyMembers & - & - & + \\ \hline
	\end{tabular}\par}

\vspace{1.0em}
\textbf{RB-6/ParentChildRelationships, FamilyMembers (RBTV liên bộ, liên quan hệ)}
\begin{itemize}
	\item \emph{Nội dung:} Mỗi thành viên (ở vai trò con) có tối đa hai quan hệ cha/mẹ loại \texttt{'biological'}.
\end{itemize}

\begin{lstlisting}[language=SQL]
	forall c in FamilyMembers,
	|{r in ParentChildRelationships | r.ChildMemberId = c.MemberId AND r.RelationshipType = 'biological'}| <= 2
\end{lstlisting}
\noindent\textbf{Bối cảnh:} \texttt{ParentChildRelationships}, \texttt{FamilyMembers}.

\vspace{0.4em}
{\centering
	\renewcommand{\arraystretch}{1.2}
	\begin{tabular}{|>{\raggedright\arraybackslash}p{0.22\textwidth}|>{\centering\arraybackslash}p{0.22\textwidth}|>{\centering\arraybackslash}p{0.22\textwidth}|>{\centering\arraybackslash}p{0.22\textwidth}|}
		\hline
		\textbf{RB-6} & \textbf{Insert} & \textbf{Delete} & \textbf{Update} \\ \hline
		ParentChildRelationships & + & - & + \\ \hline
		FamilyMembers & - & - & - \\ \hline
	\end{tabular}\par}

\vspace{1.0em}
\textbf{RB-7/MemberVerifications, FamilyMembers, UserRoles (RBTV tham chiếu liên quan hệ)}
\begin{itemize}
	\item \emph{Nội dung:} Người duyệt yêu cầu xác minh (\texttt{VerifiedBy}) phải có vai trò (\texttt{UserRoles}) trong đúng dòng họ của thành viên được xác minh.
\end{itemize}

\begin{lstlisting}[language=SQL]
	forall v in MemberVerifications (v.VerifiedBy IS NULL OR
	exists m in FamilyMembers, exists ur in UserRoles
	(m.MemberId = v.MemberId AND ur.UserId = v.VerifiedBy AND ur.FamilyId = m.FamilyId))
\end{lstlisting}
\noindent\textbf{Bối cảnh:} \texttt{MemberVerifications}, \texttt{FamilyMembers}, \texttt{UserRoles}.

\vspace{0.4em}
{\centering
	\renewcommand{\arraystretch}{1.2}
	\begin{tabular}{|>{\raggedright\arraybackslash}p{0.22\textwidth}|>{\centering\arraybackslash}p{0.22\textwidth}|>{\centering\arraybackslash}p{0.22\textwidth}|>{\centering\arraybackslash}p{0.22\textwidth}|}
		\hline
		\textbf{RB-7} & \textbf{Insert} & \textbf{Delete} & \textbf{Update} \\ \hline
		MemberVerifications & + & - & + \\ \hline
		UserRoles & - & - & + \\ \hline
		FamilyMembers & - & - & + \\ \hline
	\end{tabular}\par}

\vspace{1.0em}
\textbf{RB-8/Comments (RBTV liên thuộc tính, tự tham chiếu)}
\begin{itemize}
	\item \emph{Nội dung:} Bình luận trả lời (\texttt{ParentCommentId} khác NULL) phải thuộc cùng bài đăng (\texttt{PostId}) với bình luận cha.
\end{itemize}

\begin{lstlisting}[language=SQL]
	forall c in Comments (c.ParentCommentId IS NULL OR
	exists p in Comments (p.CommentId = c.ParentCommentId AND p.PostId = c.PostId))
\end{lstlisting}
\noindent\textbf{Bối cảnh:} \texttt{Comments}.

\vspace{0.4em}
{\centering
	\renewcommand{\arraystretch}{1.2}
	\begin{tabular}{|>{\raggedright\arraybackslash}p{0.22\textwidth}|>{\centering\arraybackslash}p{0.22\textwidth}|>{\centering\arraybackslash}p{0.22\textwidth}|>{\centering\arraybackslash}p{0.22\textwidth}|}
		\hline
		\textbf{RB-8} & \textbf{Insert} & \textbf{Delete} & \textbf{Update} \\ \hline
		Comments & + & - & + \\ \hline
	\end{tabular}\par}

\vspace{1.0em}
\textbf{RB-9/Marriages (RBTV liên bộ)}
\begin{itemize}
	\item \emph{Nội dung:} Hai vợ chồng phải là hai thành viên khác nhau; tại một thời điểm mỗi thành viên không có quá 1 cuộc hôn nhân ở trạng thái \texttt{'married'}.
\end{itemize}

\begin{lstlisting}[language=SQL]
	forall m in Marriages (m.Spouse1Id != m.Spouse2Id)
	forall x in FamilyMembers,
	|{m in Marriages | (m.Spouse1Id = x.MemberId OR m.Spouse2Id = x.MemberId) AND m.Status = 'married'}| <= 1
\end{lstlisting}
\noindent\textbf{Bối cảnh:} \texttt{Marriages}.

\vspace{0.4em}
{\centering
	\renewcommand{\arraystretch}{1.2}
	\begin{tabular}{|>{\raggedright\arraybackslash}p{0.22\textwidth}|>{\centering\arraybackslash}p{0.22\textwidth}|>{\centering\arraybackslash}p{0.22\textwidth}|>{\centering\arraybackslash}p{0.22\textwidth}|}
		\hline
		\textbf{RB-9} & \textbf{Insert} & \textbf{Delete} & \textbf{Update} \\ \hline
		Marriages & + & - & + \\ \hline
	\end{tabular}\par}

\vspace{1.4em}
\phantomsection
\addcontentsline{toc}{subsection}{3.3.3. Bảng tổng hợp các ràng buộc toàn vẹn}
{\large\bfseries 3.3.3. Bảng tổng hợp các ràng buộc toàn vẹn}\par\vspace{0.6em}

\begingroup
\renewcommand{\arraystretch}{1.25}
\setlength{\tabcolsep}{5pt}
\begin{longtable}{|>{\raggedright\arraybackslash}p{0.1\textwidth}
		|>{\raggedright\arraybackslash}p{0.22\textwidth}
		|>{\raggedright\arraybackslash}p{0.32\textwidth}
		|>{\raggedright\arraybackslash}p{0.26\textwidth}|}
	\hline
	\textbf{Mã} & \textbf{Loại RBTV} & \textbf{Quan hệ liên quan} & \textbf{Cơ chế thực thi đề xuất} \\ \hline
	\endfirsthead
	\hline
	\textbf{Mã} & \textbf{Loại RBTV} & \textbf{Quan hệ liên quan} & \textbf{Cơ chế thực thi đề xuất} \\ \hline
	\endhead
	RB-1 & Miền giá trị & Users & CHECK (Phụ lục A) \\ \hline
	RB-2 & Liên thuộc tính & Events & CHECK / Trigger \\ \hline
	RB-3 & Liên thuộc tính & FamilyMembers & CHECK / Trigger \\ \hline
	RB-4 & Liên bộ (tự tham chiếu) & ParentChildRelationships & CHECK (CK\_NotSelfParent) \\ \hline
	RB-5 & Tham chiếu liên quan hệ & EventParticipants, Events, FamilyMembers & Trigger \\ \hline
	RB-6 & Liên bộ, liên quan hệ & ParentChildRelationships, FamilyMembers & Trigger \\ \hline
	RB-7 & Tham chiếu liên quan hệ & MemberVerifications, FamilyMembers, UserRoles & Trigger \\ \hline
	RB-8 & Liên thuộc tính (tự tham chiếu) & Comments & Trigger \\ \hline
	RB-9 & Liên bộ & Marriages & CHECK + Trigger \\ \hline
\end{longtable}
\endgroup

\vspace{0.8em}
\noindent\textbf{Nhận xét:} Các ràng buộc đơn quan hệ (RB-1 $\rightarrow$ RB-4) được thực thi hiệu quả bằng \texttt{CONSTRAINT CHECK}. Các ràng buộc liên quan hệ phức tạp (RB-5 $\rightarrow$ RB-9) đòi hỏi cơ chế Trigger ở tầng cơ sở dữ liệu hoặc xử lý tại tầng ứng dụng.
